\centerline{\bf \Huge Региональная тренировка I}

\begin{flushright}
        \scriptsize{Каждый ноль за задачу добавляет тыщу задач в домашку}
\end{flushright}

\problem{1}
Попарно различные натуральные числа $a, b$ и $c$ таковы, что числа $\frac{1}{a}, \frac{1}{b}$ и $\frac{1}{c}$ образуют возрастающую арифметическую прогрессию. Докажите, что наибольший общий делитель чисел $a$ и $b$ больше 1.

\problem{2}
Дана таблица $2n\times2n$. Эльзята записала в каждую клетку таблицы либо 0, либо 1. Оказалось так, что нулей и единиц в таблице поровну. Докажите, что найдутся две строки или два столбца с одинаковым числом единиц.

\problem{3}
Все коэффициенты многочлена $P(x)$ равны либо 0 , либо 1 , причём $P(1)=25$. Может ли число 11 встретиться среди коэффициентов многочлена $P^2(x)$ хотя бы 16 раз?

\problem{4}
Внутри треугольника $A B C$ выбрана точка $X$. Точки $P, Q$ и $R$ - вторые точки пересечения прямых $X A, X B$ и $X C$ с описанной окружностью треугольника $A B C$. Пусть $U$ - точка на отрезке $X P$. Прямая, проходящая через $U$ параллельно $A B$, пересекает отрезок $B Q$ в точке $V$. Прямая, проходящая через $U$ параллельно $A C$, пересекает отрезок $C R$ в точке $S$. Докажите, что точки $Q, R, V$ и $S$ лежат на одной окружности.

\problem{5}
Дана доска размером $100 \times 100$ клеток. Мышь находится в левой верхней клетке, а кусок сыра - в правой нижней. Мышь хочет добраться до сыра. За один шаг мышь может переместиться в соседнюю по стороне клетку. На любое общее ребро двух соседних клеток можно поместить разделитель, который запрещает прямое перемещение между этими клетками. Размещение разделителей называется \textit{добрым}, если мышь всё ещё может достичь сыра после их установки. Найдите наименьшее $n$ такое, что при любом добром размещении 2023 разделителей мышь может добраться до сыра не более чем за $n$ шагов.